\section{Wave propagation in homogeneous layers}
In the previous section, a formulation for rigid-frame porous media and two formulations for poroelastic media modelling have been presented. Concluding the presentation of these motion equations, some characteristics of the wave propagation in unbounded porous and poroelastic media have been shown. The present section proposes the introduction of a geometry, that is to say boundaries or interface conditions delimiting different domains characterized each by a different type of material. Generally, the configurations studied in this work assume a monochromatic plane wave originated from a fluid half-space, i.e. air, impinging at the surface of a structure and is then scattered (reflected and transmitted) from and through each interface. This is illustrated in Fig.~\ref{fig:refl_transm_pb_schema} for the two main cases that are the so-called reflection problem where the structure is bounded on one side by a rigid-backing and the transmission problem where the incident wave is reflected and transmitted through the layer.

If the material in which the wave propagates is bounded by two parallel, free, and flat surfaces, i.e., a plate, then it is refered to as a layer. The waves supported by the medium propagate by reflecting alternatively from one surface to the other; they are said to be guided along the layer. These observations of successive reflections and refractions lead to the concept of waveguides, but their analytical descriptions often leads to writing complex system that can be difficult to solve. To determine the modes of a structure, one must seek to solve the propagation equations and the boundary conditions imposed at its boundaries. In a first section, the scattering of a plane wave by a homogeneous porous layer is investigated along with the various method available. Next, the methods for the computation of the modes of these configurations are presented. In a subsequent section, more complex geometries, those of metaporous structures, are of interest and both methods for the computation of the scattering coefficients and the modes of those configurations will be discussed.

\begin{figure}[htpb]
    \centering
    

\begin{tikzpicture}
\def\height{1.5} \def\pad{.3} \def\arrow{.7}\def\totalen{7}
\begin{scope}[xshift=-7cm, yshift=0,scale=.75]
    \fill [
        pattern={Lines[distance=2mm,angle=45,line width=0.7mm]},pattern color=black
       ] (0, \height) rectangle ++(\totalen, .5*\pad);
    \draw[top color = fluid, bottom color = white,draw=none] (0, 0) rectangle ++ (\totalen, -6*\pad);
    \draw[fill=pem, draw=none] (0, 0) rectangle ++ (\totalen, \height);

    \draw[black] (0, 0) -- ++ (\totalen, 0);
    \draw[black] (0, \height) -- ++ (\totalen, 0);
    \draw[<->] (-2*\pad, 0) -- ++ (0, \height) node[midway, left] {$h$};
    \draw[thick, ->] (0, 0) -- ++ (0, \arrow) node[right] {$x_2$};
    \draw[thick, ->] (0, 0) -- ++ (\arrow, 0) node[below] {$x_1$};
    \draw[thick, <-] (\totalen/2-2*\pad, \pad) -- ++ (0, \height-2*\pad) node[midway, left] {$b_-$};
    \draw[thick, ->] (\totalen/2+2*\pad, \pad) -- ++ (0, \height-2*\pad) node[midway, right] {$b_+$};
    \draw[thick, <-] (\totalen/2, 0) -- ++ (210:4*\pad) node[left] {$a_+$};
    \draw[thick, ->] (\totalen/2, 0) -- ++ (-30:4*\pad) node[right] {$R$};
    \draw[dashed] (\totalen/2, 0) -- ++ (0, -4*\pad);
    \draw[<->] (\totalen/2, -2*\pad) arc (90:30:-2*\pad) node[midway, below] {$\theta$};
    \node at (-1, 2.3){a)};
\end{scope}
\begin{scope}[xshift=0cm, yshift=0, scale=.75]
    \draw[top color = fluid, bottom color = white,draw=none] (0, 0) rectangle ++ (\totalen, -6*\pad);
    \draw[top color = white, bottom color = fluid,draw=none] (0, \height) rectangle ++ (\totalen, 6*\pad);
    \draw[fill=pem, draw=none] (0, 0) rectangle ++ (\totalen, \height);

    \draw[black] (0, 0) -- ++ (\totalen, 0);
    \draw[black] (0, \height) -- ++ (\totalen, 0);
    \draw[<->] (-2*\pad, 0) -- ++ (0, \height) node[midway, left] {$h$};
    \draw[thick, ->] (0, 0) -- ++ (0, \arrow) node[right] {$x_2$};
    \draw[thick, ->] (0, 0) -- ++ (\arrow, 0) node[below] {$x_1$};
    \draw[thick, <-] (\totalen/2-2*\pad, \pad) -- ++ (0, \height-2*\pad) node[midway, left] {$b_-$};
    \draw[thick, ->] (\totalen/2+2*\pad, \pad) -- ++ (0, \height-2*\pad) node[midway, right] {$b_+$};
    \draw[thick, <-] (\totalen/2, 0) -- ++ (210:4*\pad) node[left] {$a_+$};
    \draw[thick, ->] (\totalen/2, 0) -- ++ (-30:4*\pad) node[right] {$R$};
    \draw[thick, ->] (\totalen/2, \height) -- ++ (30:4*\pad) node[right] {$T$};
    \draw[dashed] (\totalen/2, 0) -- ++ (0, -4*\pad);
    \draw[<->] (\totalen/2, -2*\pad) arc (90:30:-2*\pad) node[midway, below] {$\theta$};
    \node at (-1, 2.3){b)};
\end{scope}
\end{tikzpicture}
    \caption{Diagram representing the cases of the a) reflection and b) transmission problem}
    \label{fig:refl_transm_pb_schema}
\end{figure}

\subsection{Analytical plane wave modelling}
\label{ssec:GM}

Modeling acoustic propagation in homogeneous layers by assuming plane waves is the simplest formalism in acoustics. We assume these plane waves such that they satisfy the motion equations, e.g. the ones introduced in \cref{eq:chap1_helmholtz} or \cref{eq:chap1_biot}. Introducing a simple problem, that of the scattering of an incident plate by a porous layer, we present an example of plane wave modelling for a simple system. The system is 2D with coordinates $\mbf x = (x_1, x_2)$ constituted of a porous layer with thickness $h$ and thus two interfaces of infinite extent along $x_1$, there is a fluid-porous interface at $x_1 = h$ and a porous-fluid interface at $x_1 = 0$. Both media of the air half-spaces and porous layer are modeled as fluid, thus the interfaces conditions should impose continuity of the pressure fields and normal particle displacements. The incident pressure field that satisfy the Helmholtz equation and originated from the lower interface takes the form of, 
\begin{equation}
    p_{0-} = a_+ e^{-ik_1 x_1} e^{-ik_{2,0} x_2} + R e^{-ik_1 x_1} e^{ik_{2,0} x_2}, 
\end{equation}
with $a_+$ the incident wave amplitude and $R$, the amplitude of the reflected wave. The pressure field in the porous layer is expressed as a sum of forward and backward waves,
\begin{equation}
    p_{p} = b_+ e^{-ik_1 x_1} e^{-ik_{2,eq} x_2} + b_- e^{-ik_1 x_1} e^{ik_2 (x_2-h)},
\end{equation}
and at the upper interface,
\begin{equation}
    p_{0-} = T e^{-ik_1x_1}e^{-ik_{2,0} (x_2-h)},
\end{equation}
with $T$ the transmitted wave amplitude. All wavenumbers should hold for the Snell-Descartes law that states
\begin{equation}
    k_1^2 + k_{2,0}^2 = k_0^2, \,\,\mathrm{ and }\,\, k_1^2 + k_{2,eq}^2 = k_{eq}^2. 
    \label{eq:snell}    
\end{equation}
The wavenumber should be imposed with $k_1 = k_0 \sin{\theta}, k_{2,0} = -k_0 \cos{\theta}$, with $\theta$ the angle of incidence, that governs the direction of propagation of the incident field. The displacement fields are additionally expressed using the momentum equation relating the latter with the pressure fields as $\partial_2 p = -\im \omega \rho_i u_2$, with $i=0, eq$. 

Next, the interface conditions describe the relationships between properties at the boundary between two media, including the continuity or vanishing of given fields. If a media is terminated by a surface (e.g. a rigid wall), a constraint should be applied on the fields to represent it, denoted as boundary conditions. Assuming a continuity between the two fluids in the present case with interfaces located at $x_2 = 0$ and $x_2 = h$, we obtain the following conditions,
\begin{equation}
    p^t(h) = p^l(h), \quad u_{2}^t(h) = u_2^l(h) \quad p^l(0) = p^b(0) \quad u_{2}^l(0) = u_2^b(0),
\end{equation}
with superscripts $b$ and $t$ denoting respectively the bottom and top interfaces. 

The full analytical description of a system constituted from one or several interfaces can be done by writing a linear system built from the interface conditions, which vector of unknowns are the amplitudes of each of the existing waves in the system, that is $\begin{pmatrix} b_+ & b_- & R & T \end{pmatrix}^\transp$. The forcing right-hand side corresponds to the incident wave terms. The solution of the system yields all unknown amplitudes of the system, in particular the reflection and transmission coefficients. Imposing the interface conditions leads to writing the following matrix system, 
\begin{equation}
    \begin{pmatrix}
        1 & e^{\im k_{2,eq}h}& 0 & -1 \\
        \frac{k_{2,eq}}{\rho_{eq}}& -\frac{k_{2,eq}}{\rho_{eq}} e^{\im k_{2,eq}h}&   0& \frac{k_{2,0}}{\rho_0} \\
        e^{\im k_{2,eq}h}& 1 & -1 & 0 \\
        \frac{k_{2,eq}}{\rho_{eq}} e^{\im k_{2,eq}h} &    -\frac{k_{2,eq}}{\rho_{eq}}& -\frac{k_{2,0}}{\rho_{0}}&  0 
    \end{pmatrix} \begin{pmatrix} b_+ \\ b_- \\ R \\ T \end{pmatrix} = \begin{pmatrix}
        0 \\ 0 \\ a_+ \\ \frac{k_{2,eq}}{\rho_{eq}} a_+
    \end{pmatrix}. 
    \label{eq:gm_pres_layer}
\end{equation}

From this expression, the system can simply be solved to obtain all scattered amplitudes $b_+, b_-, R$ and $T$, where $R$ and $T$ correspond respectively to the reflection and transmission coefficient. In the case of a rigidly-backed layer, the interface at $x_2=h$ is replaced by a rigid termination. For (equivalent) fluids, a hard wall or Neumann-type boundary condition is translated by a vanishing normal particle displacement field along the boundary, that is $u_2(\Gamma) = 0$. The term depending on $T$ thus vanishes, since $p_{0+}$ does not exist; this reduces the size of the system.

This method is denoted in this work as the \acrfull{gm}, since all of the available data is used in the formulation of the problem as the unknowns are all wave amplitudes. As described in the method introduced in Chapter \ref{chap:pb_direct}, a general formalism can be used for this method and any stack of any layer can be described with the GM, by putting together a matrix system that contains all interface conditions in the system, in the general form,
\begin{equation}
    \mbf M(\omega) \mbf x = \mbf f,
    \label{eq:sys_gm}
\end{equation}
where $\mbf M$ corresponds to the interface conditions involved in the system, $\mbf x$ contain the wave amplitudes and $\mbf f$ is the forcing term, the incident plane wave amplitudes. 

Similarly, a linear system can be written in the case of a \acrshort{pem} layer; the outer pressure fields remain the same but a solid and a fluid displacement field is modeled in the layer, added to the pressure field, as mentioned in Section~\ref{sec:porous_models}. The interface conditions for a poroelastic-fluid coupling are thus, 
\begin{equation}
\begin{aligned}&\sigma_{12}^{s}=0,\quad\sigma_{22}^{s}+\sigma_{22}^{f}=-p_0,&\sigma_{22}^{f}=-\phi p_0, \quad\phi u^f_{2}+(1-\phi)u^s_{2}=0.\end{aligned}\end{equation}
and for a rigid-backing, 
\begin{equation}
    u^s_1(\Gamma) = 0, \quad u^s_2(\Gamma) = 0, \quad u^f_2(\Gamma) = 0.
\end{equation}
For the case of a rigidly-backed layer, the system is entirely written in in the appendices of Ref.~\cite{weisser2016}. For the symmetric layer, the linear system in given in Ref.~\cite{boeckx2005}. Finally, the procedure to obtain the scattering coefficients is the same as previously. 

A few results for the simple cases of the scattering of a plane wave by a rigidly-backed layer and in the case of a symmetric configuration (with air coupling at both interfaces) are represented by the diagrams in Figs~\ref{fig:poroelastic_scattering_ex}a and b and results in Figs~\ref{fig:poroelastic_scattering_ex}c and d, respectively. The GM results are presented by the continuous lines. There, the layer for both porous and PEM are both $h=5$ cm thick and the incident wave propagates with a $\theta=30^o$ angle. The parameters for each material corresponds to the one used in Table~\ref{tab:poroelastic_paramsJAP} from Chapter~\ref{chap:jap}.

The first case, also referenced to as a reflection problem compares the equivalent fluid and the full Biot model response. The two curves follow the same trend but additional effects due to the skeleton motion can be observed on the continuous lines while the equivalent fluid response shows to be smoother. For all configurations, the melamine layer is acoustically transparent at lower frequency and absorption can be noticed at slightly higher frequencies. 


\subsection{Transfer Matrix Method}
The widely used method we know today as the \acrfull{tmm}, was initially developed by the works of Thomson \cite{thomson1950} and Haskell \cite{haskell1953} in order to describe elastic wave propagation in multilayer systems. This method can be applied to a range of wave physics among geophysics, quantum mechanics, electromagnetics and obviously acoustics \cite{bruneau2006, markos2008,jimenez2021} and elastic waves for ultrasonics\cite{lowe1995}. Its principle lies in the evaluation of the state vector that regroups the variables of interest in the system and evaluate them at each interface, which is possible thanks to the interface conditions. Then, the physical field amplitudes stored in the state vector are propagated from one layer to the next in the case where multiple coupled layers are modeled. The transfer matrix $\mbf T$ of a fluid layer, relating the pressure and normal displacement fields from the bottom (located at $x_2 = 0$) to the top interface (located at $x_2 = h$) of a fluid layer is thus, \cite{bruneau2006}

\begin{equation}
    \begin{pmatrix}
        p(0) \\ u_2(0)
    \end{pmatrix} = 
   \underbrace{ \begin{pmatrix}
        \cos(k_{eq}h) & \im Z_{eq} \sin(k_{eq}h) \\
        \im Z_{eq}^{-1} \sin(k_{eq}h) & \cos(k_{eq}h) \\
    \end{pmatrix}}_{\mbf T}
    \begin{pmatrix}
        p(h) \\ u_2(h)
    \end{pmatrix}, 
\end{equation}
where $k_{eq}$ is the bulk wavenumber of the wave propagating in the equivalent fluid and $Z_{eq}$ is the medium impedances. The extension to \acrshort{pem} makes of a state vector that includes the fields, $ \begin{pmatrix} \hat{\sigma}_{12} & u_2^s & u_2^t & \hat{\sigma}_{22} & p & u_1^s \end{pmatrix}^\transp$. More details are given in Chapter 11 of Allard and Atalla textbook \cite{allard2009}.

Stacks of $N$ layers that have the same material nature are easy to manipulate, the transfer matrix of each layer has the same size as the others and thus,
\begin{equation}
\mbf T_{\mathrm{global}} = \prod_{n=0}^N \mbf T_n.
\end{equation}
\acrshort{tmm} allows for an easy multiphysics coupling with interface matrices that relate the relations between the various fields at an interface which are thus non square matrices. This method was deep dived in several textbooks among which we propose Refs.~\cite{allard2009, jimenez2021} for additional details. In the last decades, there has been some works focused in the extension of this method to stabilize it and use for different applications and various formulations have been introduced like the Global Matrix Method \cite{knopoff1964} which assembles the transfer matrices of all layers into a global equation system. The Stiffness Matrix Method has also been investigated and uses recursively applied compliance matrices to obtain a computationally stable formulation \cite{rokhlin2002}. A different recursive approach has also been proposed by Dazel \etal \cite{dazel2013} and introduces an Information Vector used to stabilize the method. Variations of the TMM also allows for the study of 2D system by combining transfer matrices for series elements and admittance matrices for parallel elements. This method has been used to model a \acrfull{hr} embedded in an equivalent fluid porous layer \cite{doutres2015}. 

\begin{figure}[htpb]
    \centering 
    \includegraphics{poro_elastic_scattering.pdf}
    \caption{Scattering coefficients from a poroelastic layer (continuous) and equivalent fluid layer (dashed) computed using the GM for the a) reflection, b) transmission problem. Results computed with the TMM are also shown for the absorption coefficient for the equivalent fluid (\smallcirclemarker) and for the poroelastic layer (\smalltrianglemarker)}
    \label{fig:poroelastic_scattering_ex}
\end{figure}

While very versatile, it remains unstable in limit cases at high frequencies or for large thickness of the layers. While the formulation is mathematically exact, the evaluation of the system involves involves finite-precision calculation which leads the very sensitive exponential terms to diverge in the previously specified regimes. From the evaluation of the transfer matrix of the global system, the scattering matrix can be obtained and thus the reflection and transmission coefficients by simply manipulating the expressions for the fields contained in the state vector with the scattering coefficients. 

Considering a rigidly-backed layer, and thus $u_2(h) = 0$, we obtain a reflection coefficient,
\begin{equation}
    R_b =\frac{T_{11}-Z_{eq}T_{21}}{T_{11}+Z_{eq}T_{21}}.
\end{equation}

In the case of the symmetric layer, the reflection and transmission coefficients are determined from the scattering matrix as well, 
\begin{equation}
    T_s = \frac{e^{-ikL}}{T_{11}+T_{12}/Z_{eq}+Z_{eq}T_{21}+T_{22}}, \quad
    R_s = \frac{T_{12}/Z_{eq}-Z_{eq}T_{21}}{T_{11}+T_{12}/Z_{eq}+Z_{eq}T_{21}+T_{22}}.
\end{equation}
which is only valid for reciprocal and linear assumptions. The absorption coefficient for the same two configurations are calculated as $\alpha_b = 1 - |R_b^2|$ and $\alpha_s = 1 - |R_s^2| - |T_s^2|$ and are shown in Fig.~\ref{fig:poroelastic_scattering_ex}with markers. It can be seen that the same results (up to the respective numerical precision of each method) are obtained. No absorption peak is observed in the results since the frequency range is way below the quarter-wavelength resonance (which should lie around 1.5 kHz for these 5 cm-thick homogeneous layers). The equivalent fluid and the full Biot theory results follow generally the same trend for the results of the rigidly-backed layer in Fig.~\ref{fig:poroelastic_scattering_ex}a, except for a peak around 200 Hz shown with the poroelastic model. This is known as the Biot peak and corresponds to an elastic resonance, a strong vibration of the skeleton leading to increased structural dissipations, hence the reflection coefficient dip. \cite{zielinski2011, weisser2016}. The results in the transmission case in Fig.~\ref{fig:poroelastic_scattering_ex}b present a generally lower absorption, bounded at $\alpha = 0.5$. Here, the equivalent fluid and poroelastic respective response are significantly different: at low frequencies, the Biot theory behaves as expected with the transmission close to 1 and the reflection close to 0. The equivalent fluid reflection coefficient is similar but its transmission lies around 0.6 for lower frequencies, leading to a rather important absorption coefficient at zero frequency. While not representing a physical and observable behaviour of the system, this effect is a feature of the JCA model due to the low-frequency limit of the real part of the effective density. A correction term has been proposed by Pride \etal \cite{pride1993} to improve this low-frequency limit \cite{allard2009}.

\section{Dispersion of guided waves in layers}
\label{sec:chap1_guidedwaves}
In this section, we take an interest in similar configurations as previously, removing the excitation or forcing term in the equations.  Solving the appropriate set of equations for a given set of parameters leads to the modes of the free system, that is the solution of the characteristic equation leading a specific behaviour of the wave propagation in the structure for the said parameters. The form of the modes of wave propagation in layer is discussed and various methods for their calculation of are presented. The existence of an interface in the geometry leads to the propagation of the so-called surface waves \cite{auld1973} where the interface itself acts as the support for propagation. 

\subsection{Example of a planar air waveguide}
A dispersion relation can be illustrated using the example of an elementary planar waveguide. A fluid layer of infinite extent in one direction, with thickness $h$, and Neumann-type boundary conditions at each end $x_2 = 0$ and $x_2 = h$. This configuration is the simplest type of waveguide where waves are guided by successive reflection onto each wall. This leads to modes of the system in the form $k_{2,n} = n\pi/h, \forall n \in \mathbb N$ \cite{bruneau2006,royer2021}. Finally, the dispersion relation writes as,
\begin{equation}
    k_{1,n}(\omega) = \sqrt{\left(\frac{\omega}{c_0}\right)^2 - \left(\frac{n \pi}{h}\right)^2},
\end{equation}
with $c_0$ the velocity in air. The solutions to this equation for a range of $n=0$ to $n=4$ is shown in Fig.~\ref{fig:disp_air}. There, the wavenumber appears as either purely real or purely imaginary, indicating respectively the propagative and evanescent modes of this simple geometry. The mode with order 0 is the one that exists in unbounded fluid materials, with its slope corresponding to the fluid velocity $c_0$. The next orders all present a cut-off frequency below which the modes become evanescent.

\begin{figure}[htpb]
    \centering
    \includegraphics{disp_airlayer.pdf}
    \caption{Adimensional dispersion relation of an air waveguide with thickness $h=5$ cm. a) Real part of the wavenumber and b) Imaginary part of the wavenumber}
    \label{fig:disp_air}
\end{figure}

Different types of waves propagate along the interface. An example can be water waves that follow the interface between air and water; most of the energy is localized at the interface. A similar behaviour exists for seismic waves where the interface of a solid substrate supports the propagation of Rayleigh (longitudinal motion) and Love waves (transverse motion) which both show an energy strongly decaying with depth. Scholte waves are waves that propagate along the interface between a fluid and a solid medium. In fluid layers, the propagating waves are similar to studying a shear horizontal (SH) wave in an elastic layer. One mode propagates at zero frequency, while countless other modes propagate with a cut-off frequency that corresponds to the successive quarter-wavelength resonances, essentially depending on the layer thickness and the nature of the boundary/interface conditions. 

From the previously obtained linear systems in Section~\ref{ssec:GM} for the configurations, and removing the incident wave from the system, i.e. the forcing (right-hand side) term, an homogeneous linear system is obtained, in the general form $\mathcal L(\mbf k, \omega) \mbf s = 0$, where the operator $\mathcal L$ governs wave propagation in the structure, $\mbf k = (k_1, k_2)^\transp$ is the wavenumber vector in the 2D geometry and $\mbf s$ a vector containing the amplitude of all the partial waves. In order to obtain nontrivial solutions for $\mbf s$, the system should satisfy a vanishing determinant for the linear system, $\mathrm{det}(\mbf L (\mbf k, \omega)) = 0$. The parameters satisfying this so-called \emph{dispersion equation} can be sought, allowing to identify the modes of the structure. These equations can be solved for various parameter spaces depending on the perpective adopted on the problem \cite{ducasse2022, deschamps2025}: solving for frequencies while imposing the wavenumber or the other way around. The first one is useful when interested in the time-dependant effects the structure has on the propagation. In this work, we are interested in the losses generated by the structure and the material and as such frequency is fixed in order to solve complex wavenumbers. The solutions are thus able to indicate on rather the spatially induced dispersion occuring in the system.

\subsection{A bestiary of wave solutions in dissipative and open waveguides}
\label{sec:bestiary}

In this section, an overview of the possible sets of solutions for the guided waves in waveguide is presented. These sets of solutions are determined by the direction of the propagation, the attenuation and the nature of the leaky wave if it exists. This \emph{bestiary} is split into 3 categories. Considering a layer and half-space system, there exists three wavenumber components from Eq.~(\ref{eq:snell}). The form of solutions are assumed as 
\begin{equation}
   p_l \propto e^{-ik_1 x_1} e^{\pm ik_{2,eq} x_2}, \quad p_0 \propto e^{-ik_1 x_1} e^{\pm ik_{2,0} x_2}, 
\end{equation}
with $k_1 = +|k_1'| - i|k_1''|$ and $k_{20} = \pm |k_{20}'| \pm i|k_{20}''|$. The longitudinal wavenumber is the same in both layer and half-space domain as a consequence of the Snell-Descartes law and the continuity equations: there exists in total three dependant wavenumber components in this system, $k_1$, $k_{2,0}$ and $k_{2,eq}$ among which determining two allows to compute all,
\begin{equation}
    k_{2,0} = \pm \sqrt{k_0^2 - k_1^2}, \quad k_{2,eq} = \pm \sqrt{k_{eq}^2 - k_1^2}. 
\end{equation}
An important relation between these wavenumbers can be obtained by taking the imaginary part of Eq.~\ref{eq:snell},
\begin{equation}
    k_1'k_1'' = -k_{20}' k_{20}''. 
    \label{eq:treyssedeexpr_disp}
\end{equation} 
A set of physically meaningful types of solutions should satisfy this equation, \cite{treyssede2014}; the following section will provide a detailed discussion of this issue.

Because of the symmetrical nature of the presented configurations, and the considerations of unbounded layers in the $x_1$ direction, a plane of solutions corresponding to $k_1' \geq 0$ can be selected, as the right-going solution is perflectly equivalent solutions to the left-going solutions.  
\begin{figure}[ht]
    \centering
    \begin{tikzpicture}
    \def\arrowlen{0.7}
    \def\len{2}
    \def\hei{1}
    \def\pad{0.4}
    \begin{scope}[shift={(0,0)}]
        \draw node at (-.5, 1.5) {c)};
        \node at (\len/10, \hei+2*\pad) {I};
        \draw (0, 0) -- ++ (\len, 0);
        \draw (0, \hei) -- ++ (\len, 0);
        \draw[blue!60!black, thick,samples=100,smooth,variable=\t,domain=0:\len/2]
                plot(\t+\len/5,{(\hei-\pad)/2*exp(-\t/.5)*cos(800*\t)+\hei/2});
        \draw[thick, Latex-] (6.5*\len/10, \hei+\pad) -- ++ (40:.6*\arrowlen);
        \draw[thick, Latex-] (5*\len/10, \hei+\pad) -- ++ (40:.8*\arrowlen);
        \draw[thick, Latex-] (3.5*\len/10, \hei+\pad) -- ++ (40:1.*\arrowlen);
        \draw[->] (0, 0) -- ++ (\hei/2, 0) node[below right] {$x_1$};
        \draw[->] (0, 0) -- ++ (0, \hei/2) node[above left] {$x_2$};
    \end{scope}
\end{tikzpicture}
    \caption{Classes of solutions for the different types of propagating waves in open non-dissipative waveguides a) perfectly guided solution with the amplitude of the field along $x_2$ displayed in green: the solution for the field in the surrounding medium is purely evanescent b) leaky growing waves with an attenuated solution in the layer and growing outgoing solution in the surrounding domain, c) backward leaky attenuated waves corresponding to an ingoing solution, d) outgoing growing wave solution (improper) and e) incoming growing waves (improper).}
    \label{fig:chap1_classification}
\end{figure}

In \textbf{conservative and closed systems}, the study is restricted to only the two wavenumber components in the layer, since no half-space is accounted. The total acoustic energy is thus confined to the structure enclosed by its boundaries.  A type of wave solution exists as \textit{perfectly guided}, this propagating wave has $k_1' > 0$ with no attenuation $k_1'' = 0$. Solutions include \textit{purely evanescent} wave as well $k_1'' > 0$ which often exists and continuates a propagating mode below its cut-off frequency as in simple waveguides configurations as presented in Section~\ref{sec:chap1_guidedwaves}.

\textbf{Open systems} support the existence of the previous types of modes with additional phenomena. Since part of the energy is susceptible to radiate out from the structure, the energy is no longer confined to the layer. Some of the previously purely guided modes can thus become attenuated, $k_1'' < 0$. Waves are still \textit{perfectly guided} with conditions $k_{20}' = 0$ with evanescent waves in the fluid $k_{20}'' < 0$, meaning that waves are confined to the proximity of the layer interface as shown in Fig.~\ref{fig:chap1_classification} with the amplitude along the $x_1$ direction in blue and the amplitude along the $x_2$ direction. The additional condition $k_1' = 0$ correspond to \textit{trapped modes}, which are a set of discrete eigenvalues of the system. They can exist in waveguides where the system is confined in one direction at least, the energy can thus be localized around some obstacle or defect by the evanescent waves \cite{pagneux2013, roitberg1998}. In open systems, they correspond to  waves that are evanescent at infinity. However, in open waveguides, the acoustic energy is usually able to radiate outside of the waveguide, leading to two type of waves as \textit{leaky waves} in the case $k_{20}' > 0$ and \textit{incoming waves} (also called backward leaky) with $k_{20}' < 0$. These solutions are respectively exponentially growing as it radiates $k_{20}'' > 0$ \cite{goldstone1959,monticone2015} (see the leaky growing wave represented in Fig.~\ref{fig:chap1_classification}b and incoming growing wave in Fig.~\ref{fig:chap1_classification}e) or attenuated $k_{20}'' < 0$ (see the leaky attenuated wave represented in Fig.~\ref{fig:chap1_classification}c and incoming attenuated wave in Fig.~\ref{fig:chap1_classification}d). These at first counterintuitive results ensue from the requirements that the wave should be propagating $k_1' > 0$ and attenuated $k_1'' < 0$ along the layer, and satisfying Eq.~\ref{eq:treyssedeexpr_disp}. The book of Collin \cite{collin1991} can be refered to for fundamental discussions about the location of these solutions in the complex plane. We also mention Section 3.2 from Kiefer's thesis \cite{kiefer2022} which is dedicated to the form of solutions for quasi-guided waves in the case of an elastic layer coupled to a semi-infinite fluid half-space. 

The addition of \textbf{dissipation in open systems} leads to the mixing of several loss mechanisms: the leakage of waves and the dissipations within the guide (which can be of whatever nature). In this case, the separation between incoming and leaky waves is not possible anymore. This particular aspect has been very scarcely tackled in the literature: an unpublished article is interested in an elastic waveguide \cite{kiefer2020} while in optics, a paper discusses the waves guided in nanospheres with a complex permittivity \cite{campione2011}. There, the modes are sorted as \textit{proper} and \textit{improper} solutions. This complex case is covered in a dedicated section \ref{sec:homog_layer} in chapter \ref{sec:prb}.

In the next parts of this chapter, a discussion on available methods to compute the complex dispersion relation for guided waves in single or stacks of homogeneous layers is proposed. There are semi-analytical methods that can be used, Section \ref{ssec:semi-analyt-layers} or numerical methods that involve some sort of discretization in Section \ref{ssec:numerics-layers}.

\subsection{Modes of porous layers computed with semi-analytical tools}
\label{ssec:semi-analyt-layers}
The literature on the dispersion relation of porous/poroelastic-based structures is scarce; Groby \etal \cite{groby2008} have shown the complex dispersion relation in a homogeneous equivalent fluid layer coupled to air while Boeckx \etal\cite{boeckx2005} have computed the modes propagating in a poroelastic waveguide coupled at each top and bottom interface with an outgoing plane wave propagating in an air medium. This has been done using an analytical description of the system by expressing the interface conditions in a linear system (which we call \acrshort{gm} in this work). Both of these work have used a root-finding procedure in order to solve the modes of the structure. Another work investigating the modes of a heavy fluid saturated homogeneous poroelastic layer solves the dispersion relation of the guided waves by means of an eigenvalue problem linked to the scattering matrix\cite{belloncle2003}.

Using Eq.~(\ref{eq:sys_gm}), the non-trivial solution of the problem requires that $\mathrm{det}\left(\mbf M(k_1, \omega)\right) = 0$. The determinant of the matrix from Eq.~(\ref{eq:gm_pres_layer}) can be further expressed as,
\begin{equation}
    2 \frac{k_{2,eq} k_{2,0}}{\rho_{eq} \rho_{0}} \text{cos}(k_{2,eq} h) + \im \left(\frac{k_{2,eq}^2}{\rho_{eq}^2} + \frac{k_{2,0}^2}{\rho_{0}^2}\right) \text{sin}(k_{2,eq} h) = 0.
    \label{eq:fluideq_sym_disp}
\end{equation}
This expression can not be solved analytically. Solving wavenumbers can be done through the writing and solving of a (possibly generalized) eigenvalue problem or by root-finding of solutions in the case where writing the first one appears impossible. From this expression, one can use a root-finding procedure to solve for either the term $k_{2,eq}$ or $k_1$ to obtain the dispersion relation by varying the frequency. Results are shown in \cref{fig:porous_disp}.

\begin{figure}[htpb]
    \includegraphics{disp_poro.pdf}
    \caption{Dispersion relation of a porous layer loaded by an air half-space on both sides. a) 3D view with $k_0$ (\greendashedline) and $k_{eq}$ (\pinkdashedline), b) $k_1'$ plane , c) $k_1''$ plane. d-f) Same views in the case of the poroelastic layer with $k_p$ (\greendashedline), $k_a$ (\pinkdashedline) and $k_s$ (\yelodashedline).}
    \label{fig:porous_disp}
\end{figure}

We reference to Appendix~\ref{sec:muller} for extensive details on the root-finding procedure based on the Müller method. These results show an infinite amount of modes that have increasing attenuation, only the first four are shown in Fig.~\ref{fig:porous_disp}a to Fig.~\ref{fig:porous_disp}c. A branch corresponding to $k_{eq}$ is found by the root-finding procedure, as setting $k_1 = k_{eq}$ in Eq.~\ref{eq:fluideq_sym_disp} indeed leads to a trivial solution. Only one propagating mode exists in the frequency range shown (one branch lies between the lines $k_0$ and $k_{eq}$) as is the expected behaviour of a lossy fluid layer.

A similar configuration, with a layer modeled using the full Biot theory is shown in Fig.~\ref{fig:porous_disp}d to Fig.~\ref{fig:porous_disp}f. Results are significantly different as the dispersion relation presents a large amount of branches. Alike the equivalent fluid results each bulk wavenumber constitutes a trivial solution of the characteristic equation. Since the elastic motion of the layer is modeled, quasi-Lamb modes are obtained in this case, which are all complex because of the viscoelastic and viscothermal dissipations in the poroelastic layer. In this case, the modes are not simply separated into evanescent and propagating modes. When their attenuation is low, they are quasi-guided: corresponding to propagating solution with some attenuation. These results are comparable to those of Boeckx \etal \cite{boeckx2005} and to those presented in Section~\ref{app:fluid_coupled_poroelastic}.

Generally, characteristic equations as Eq.~\ref{eq:fluideq_sym_disp} can only be obtained for very simple cases as is the case here, for a single fluid layer. The characteristic equation for Lamb modes is also well known \cite{royer2021}. However, more complicated geometry requires to solve the root-finding iterative process by using directly a numerically-computed matrix determinant which is computationally very intensive, or in many case, it is outright not possible due to the analytical complexity of the description of the system. One of the main downside of these methods is that they require multiple initial guesses to construct an interpolating function.There is no assurance of finding a solution in a given location: depending on the solution space smoothness and shape, the iterative process might get stuck in local minima, or can simply not converge in some situations. In cases where a complete dispersion relation is complicated to be constructed, more complex numerical method could thus be required; a discretization of the geometry and an interpolation of the physical fields is done as will be discussed in the upcoming section.

\subsection{Numerics for modes computation}
\label{ssec:numerics-layers}
Several numerical methods have been investigated in the last decades with the goal of computing complex modes in homogeneous layers. Guided waves are of interest in a wide range of geometries and we focus here on the case of unbounded single (or stack of) layers with two infinite interfaces as tackled in the previous sections. 

Even if a semi-analytical description of the system is possible, it can be advantageous to be able to write the problem differently, as root-finding routines often lead to challenges as discussed in earlier sections. With that in mind, several numerical methods have been introduced with the same framework: a 1D discretization of the layer along the thickness and assuming harmonic motion along the longitudinal direction.

\subsubsection{Modeling of the semi-infinite radiating space}
There is a discussion to be made on the relevance of the strategy to use for the radiating medium modelling. In the literature, two have been presented: either the coupling can be imposed locally at the interface condition, or the domain of the surrounding medium is fully modeled, accounting for some termination condition satisfying a Sommerfeld radiation condition, by using a \acrfull{pml} or some analogous absorbing layer. The use of artificial absorbing layers and the discretization at the end leads to the apparition of a continuum of modes of the cavities \cite{gallezot2017} which requires to be sorted out. A strategy based on a local coupling between the structure and the radiating medium leads to impose a form of solution for the radiation, usually as plane waves\cite{kiefer2019}. An alternative for local coupling has been investigated by associating a \acrfull{bem} formulation for the interface\cite{mazzotti2013}. 

We argue that the advantages of the imposition of a local coupling is two-fold: it requires less degrees of freedom for any sort of discretization since the interface condition describes the full motion of the radiated waves. For dispersion problems, this is particularly interesting since leaky mode solutions are thus directly obtained, without needing to filter out all radiation modes linked to the additional cavity and the resonances of the PML. Also, as explained in Section~\ref{sec:bestiary}, leaky waves are characteristic solutions of open waveguides that show exponential growth as they radiate away from the layer, which is especially ill-suited for implementing PMLs for this kind of system, although it has been achieved\cite{zuo2017, gallezot2017}. We again refer to the thesis of Kiefer \cite{kiefer2022}, especially section 3.1, for a in-depth discussion of modeling the full infinite medium against only an interface condition coupling with the layer.

\subsubsection{Numerical methods used in the literature}
Spectral methods have been around for quite some time and have been used for many numerical applications. We focus for this work on \acrfull{scm} \cite{trefethen2000}, which has been especially used for the computation of guided wave modes in elastic layers. One of the first paper that introduced this method for the computation of complex dispersion relations is that of Adamou \etal \cite{adamou2004}, while other have focused on multilayer anisotropic structures \cite{quintanilla2015} or cylinders\cite{karpfinger2008}. These references propose an implementation for closed systems only. We refer to the introduction of Chapter~\ref{chap:jap} for an extensive bibliography on \acrshort{scm} applications. 

This method is useful for the discretization of 1D geometries as is the object of the present section. In the case of simple 2D geometries, a 1D discretization along the transverse direction associated to a harmonic motion in the longitudinal direction. Even if the use of finite elements seems farfetched, this method could be relevant in contexts where boundary effects or guide with varying cross-sections are studied. However, a variation of finite elements based on a 1D discretization of the waveguide is entitled \acrfull{safe} and is akin to \acrshort{scm} which we have described in the previous paragraph. The 1D discretization is done along the thickness of the layer by regular nodes, and the fields are interpolated using finite elements shape functions. Examples of \acrfull{safe} use cases are very similar to \acrshort{scm} and include. A detailed description of the calculation of Lamb modes using a simple SAFE implementation using high-order Lobatto polynomials is proposed in Appendix.~\ref{app:safe_lamb}. There, its performance is also compared to that of SCM. The particular advantage of \acrshort{scm} is the very quick convergence allowed by the interpolation using Chebyshev polynomials, called \emph{spectral convergence} in the literature\cite{trefethen2000,canuto2006,gheorghiu2014}. As it was shown in the aforementioned appendix, we saw that it outperforms higher-order \acrshort{safe}. Thus, we will focus on this method for the modeling of multilayer homogeneous structures. 

Open elastic layers accounting for surrounding air media that couples the elastic layer have been first presented by Kiefer \etal \cite{kiefer2019} where a \acrshort{scm} was used to discretize the elastic layer. Additionally, a local coupling is modeled with a discretized interface conditions. The system of equations is then formulated as a \acrshort{gep} and can be solved with traditional eigenvalue solvers. This work forms the basis of the work done in Chapter~\ref{chap:jap}. Other implementations of the same configuration include a discretization of the fluid half-space by means of a mapping allowing the numerical solution to increase exponentially at infinity, as is a characteristic of leaky waves\cite{georgiades2022}. 

With \acrfull{safe}, which has been been widely used for similar configurations as well have computed modes of immersed waveguides \cite{zuo2017}, leaky Lamb modes using PMLs and a coupling between \acrshort{safe} and \acrshort{bem} \cite{mazzotti2013}. A detailed overview on the literature dealing with this topic is proposed in Appendix~\ref{app:safe_lamb}. Ducasse and Deschamps \cite{ducasse2022} have also formulated an eigenvalue problem by discretizing the transverse direction of the layer using finite differences.

% \begin{figure}[h]
%     \centering
%     % \begin{tikzpicture}
    \def\n{7}; \def\nodespan{0.7}; \def\xarrowlen{1}; 
    \def\layerlen{6}; \def\arrowspacing{1};
    \begin{scope}[shift={(-3,0)}]
        
        \coordinate (repere) at (-1.2*\arrowspacing, -3*\nodespan);
        % Draw the x-axis
        % geometry drawing
        \fill[porous] (0, -1*\nodespan) 
                rectangle ++ (\layerlen, 2*\nodespan);
        % Domains legends 
        \draw[thick, black] (0, -1*\nodespan) node[left] {$\Gamma_n$}
                -- (\layerlen, -1*\nodespan) ;
        \draw[thick, black] (0, 1*\nodespan) node[left] {$\Gamma_{n+1}$}
                -- (\layerlen, 1*\nodespan) ;
        \draw[thick, black] (0, -1*\nodespan) 
                -- (\layerlen, -1*\nodespan) node[midway, below] {$\vdots$};
        \draw[thick, black] (0, 1*\nodespan) 
                -- (\layerlen, 1*\nodespan) node[midway, above] {$\vdots$};
    \end{scope}
    \begin{scope}[shift={(0,0)}]
        \coordinate (repere) at (-1.2*\arrowspacing, -3*\nodespan);
        % Draw the x-axis
        % geometry drawing
        \fill[porous] (0, -1*\nodespan) 
                rectangle ++ (\layerlen, 2*\nodespan);
        % Domains legends 
        \draw[thick, black] (0, -1*\nodespan) node[left] {$\Gamma_n$}
                -- (\layerlen, -1*\nodespan) ;
        \draw[thick, black] (0, 1*\nodespan) node[left] {$\Gamma_{n+1}$}
                -- (\layerlen, 1*\nodespan) ;
        \draw[thick, black] (0, -1*\nodespan) 
                -- (\layerlen, -1*\nodespan) node[midway, below] {$\vdots$};
        \draw[thick, black] (0, 1*\nodespan) 
                -- (\layerlen, 1*\nodespan) node[midway, above] {$\vdots$};
        % Draw Chebyshev nodes
        \foreach \i in {0,...,\n} {
            \pgfmathsetmacro\x{\nodespan*cos((\i+0.5)*180/(\n+1))}
            \draw[black, fill=black] (3/4*\layerlen, \x) circle (0.8pt);
            };
            % Draw other elements
            \draw[Latex-] (3/4*\layerlen, .38) 
            to[bend left] ++ (-.4, .9) node[above] {node $x_i $};
            \draw[-Latex] (\layerlen*0.2, -0.3) 
            -- ++ (1, 0) node[midway, above] {$ k_1$};
            \draw[Circle-] (\layerlen-0.3, 0) 
            -- ++ (\arrowspacing, 0) node[right] {$\Omega_n$};
    \end{scope}
\end{tikzpicture}
% GM: represent partial waves in each domain
% tmm: write SV 
% scm: cheb nodes 
% safe: equispace nodes
% fem: periodic cell with nodes, draw a mesh
%     \caption{Discretization for each method}
% \end{figure}

